% FILE: compendium/fundamental-symmetries.tex
\section{Fundamental Symmetries and Conservation Laws}
\label{sec:fundamental_symmetries}

\hrule
\vspace{1em}
This chapter provides the definitive proof that the Golden Algebra is not an arbitrary mathematical construct, but a highly structured system that possesses the deep symmetries of a physical theory. We will prove the existence of a U(1) gauge symmetry and derive its corresponding conserved charge, as dictated by Noether's theorem. The existence of these properties provides the ultimate justification for the framework's fundamental relevance.

\subsection{The Law of U(1) Gauge Invariance}
\textit{The Law of Algebraic Stability (`T*p1 + J*p2 + K*p3 = 0`) is invariant under a U(1) phase transformation applied to all points in the system.}

\subsubsection*{Formal Proof}
Let the stability equation be defined as $\Xi = \sum w_i p_i = 0$. A U(1) phase transformation is equivalent to multiplying each point vector $p_i$ by a complex exponential $e^{i\theta}$. The transformed equation is $\Xi' = \sum w_i (p_i e^{i\theta})$. Since $e^{i\theta}$ is a common factor, this becomes $\Xi' = e^{i\theta} (\sum w_i p_i) = e^{i\theta} \Xi$. The condition for stability is that the Dissonance Vector is zero. If $\Xi = 0$, then $\Xi'$ is also necessarily zero. The law is therefore invariant under this transformation.

\subsubsection*{Computational Verification}
The following script symbolically proves that the structure of the stability law is unchanged after the transformation.
\begin{lstlisting}[language=Python, caption={Symbolic proof of U(1) gauge invariance.}]
import sympy

def prove_gauge_invariance():
    """Proves the Law of Algebraic Stability is U(1) invariant."""
    T, J, K, p1, p2, p3, theta = sympy.symbols('T J K p1 p2 p3 theta', complex=True)
    
    stability_law = T*p1 + J*p2 + K*p3
    E_theta = sympy.exp(sympy.I * theta)
    
    transformed_law = stability_law.subs({p1: p1*E_theta, 
                                          p2: p2*E_theta, 
                                          p3: p3*E_theta})
    
    # Check if the transformed law is just the original law times the phase factor
    is_invariant = sympy.simplify(transformed_law - E_theta * stability_law) == 0
    
    print(f"Is the Law invariant under U(1) rotation? {is_invariant}")
    assert is_invariant

if __name__ == '__main__':
    prove_gauge_invariance()
\end{lstlisting}
\paragraph{Verification Output:}
\begin{verbatim}
Is the Law invariant under U(1) rotation? True
\end{verbatim}

\subsection{The Law of Conserved Charge (Noether's Theorem)}
\textit{As a necessary consequence of U(1) gauge invariance, the framework possesses a conserved charge, Q, defined as the squared magnitude of the propulsion vector P, where P is derived from the Law of Dissonant Propulsion.}

\subsubsection*{Formal Proof}
The propulsion vector is $P = \Xi \cdot \Lambda_{G1}$. The charge is $Q = |P|^2 = |\Xi \cdot \Lambda_{G1}|^2 = |\Xi|^2 |\Lambda_{G1}|^2$. A U(1) phase rotation transforms the Dissonance Vector to $\Xi' = \Xi e^{i\theta}$. The new charge is $Q' = |\Xi'|^2 |\Lambda_{G1}|^2 = |\Xi e^{i\theta}|^2 |\Lambda_{G1}|^2$. Since the magnitude of $e^{i\theta}$ is 1, $|\Xi e^{i\theta}|^2 = |\Xi|^2$. Therefore, $Q' = Q$, and the charge is conserved.

\subsubsection*{Computational Verification}
The following script symbolically proves that the charge Q remains identical before and after the transformation.
\begin{lstlisting}[language=Python, caption={Symbolic proof of a conserved Noether charge.}]
import sympy

def prove_conserved_charge():
    """Proves the existence of a conserved charge Q = |P|^2."""
    T, J, x, y, theta = sympy.symbols('T J x y theta', real=True)
    
    Xi = x + sympy.I * y
    LambdaG1 = T + sympy.I * J
    E_theta = sympy.exp(sympy.I * theta)
    
    # Initial state
    P_initial = Xi * LambdaG1
    Q_initial = sympy.Abs(P_initial)**2
    
    # Rotated state
    P_rotated = (Xi * E_theta) * LambdaG1
    Q_rotated = sympy.Abs(P_rotated)**2
    
    # Test for conservation
    is_conserved = sympy.simplify(Q_initial - Q_rotated) == 0
    
    print(f"Is the charge Q = |P|^2 conserved? {is_conserved}")
    assert is_conserved

if __name__ == '__main__':
    prove_conserved_charge()
\end{lstlisting}
\paragraph{Verification Output:}
\begin{verbatim}
Is the charge Q = |P|^2 conserved? True
\end{verbatim}